% Source: hierarchical_softmax.tex
\begin{notechapter}
% ---- chapter-local definitions ----
\tikzset{
    my colour/.style={draw=#1, fill=#1!50},
    my leaf/.style={my colour=cyan,},
  }
\forestset{
  every leaf node/.style={
    if n children=0{#1}{}
  }
}
% -----------------------------------
\chapter{Hierarchical Softmax and Huffman Trees}\label{ch:hierarchical-softmax}

\section{Introduction}
Hierarchical Softmax is a technique used in natural language
processing to efficiently compute probabilities for a large
number of outcomes, such as predicting the next word in a
sentence. It organizes outcomes into a hierarchical tree
structure, reducing computational complexity. This document
provides a detailed explanation of Hierarchical Softmax,
Huffman trees, and their application in Word2Vec.

\section{Hierarchical Softmax}
\subsection{Motivation}
In traditional Softmax, computing probabilities for a large
vocabulary (e.g., 100,000 words) is computationally expensive
because it requires calculating scores for every word and
normalizing them. Hierarchical Softmax addresses this by
organizing words into a binary tree, reducing the complexity
from \( O(N) \) to \( O(\log N) \).

\subsection{Tree Structure}
The vocabulary is organized into a binary tree:
\begin{itemize}
    \item Each \textbf{leaf node} represents a word.
    \item Each \textbf{internal node} represents a decision
          point (e.g., left or right).
\end{itemize}

\subsection{Probability Calculation}
The probability of a word \( w \) given a context vector
\( w_i \) is computed by traversing the tree from the root
to the leaf node corresponding to \( w \). At each internal
node, a binary classifier decides whether to go left or
right. The probability of \( w \) is the product of the
probabilities along the path.

\subsection{Example}
For a vocabulary [cat, dog, apple, banana], the tree might
look like this:

\begin{verbatim}
        Root
       /    \
     A       B
    / \     / \
 cat dog apple banana
\end{verbatim}

To compute \( P(\text{cat} \,|\, w_i) \):
\begin{enumerate}
    \item Start at Root.
    \item Go left to A with \( P(\text{left} \,|\, \text{Root}, w_i) \).
    \item Go left to "cat" with \( P(\text{left} \,|\, A, w_i) \).
    \item \( P(\text{cat} \,|\, w_i) = P(\text{left} \,|\, \text{Root}, w_i)
          \times P(\text{left} \,|\, A, w_i) \).
\end{enumerate}

\section{Huffman Trees}
\subsection{What is a Huffman Tree?}
A Huffman tree is a binary tree used for efficient encoding.
It minimizes the average path length by placing frequent
elements closer to the root. In Word2Vec, it is used to
organize the vocabulary for Hierarchical Softmax.

\subsection{Construction Algorithm}
The Huffman tree is constructed using a greedy algorithm:
\begin{enumerate}
    \item Create a leaf node for each word, weighted by its frequency.
    \item Repeatedly combine the two nodes with the smallest
          frequencies into a new internal node.
    \item Continue until all nodes are combined into a single tree.
\end{enumerate}

\subsection{Example}
For words [cat (10), dog (8), apple (6), banana (5), 
elephant (4), zebra (3), giraffe (2), lion (1)], the Huffman
tree is constructed as follows:\\
\hspace{2pt}\\
\begin{forest}
for tree={
    circle,
    draw=blue,
    thick,
    fill=blue!10, 
    text=red,
    s sep=25pt, l sep=10pt,
    every leaf node={my leaf},
    minimum size=6.5em,}
[Root (39)
    [Node (17)
        [dog (8)]
        [Node (9)
            [banana (5)]
            [elephant (4)]
        ]
    ]
    [Node (22)
        [cat (10)]
        [Node (12)
            [apple (6)]
            [Node (6)
                [zebra (3)]
                [Node (3)
                    [giraffe (2)]
                    [lion (1)]
                ]
            ]
        ]
    ]
]
\end{forest}

\subsection{Why Use a Huffman Tree?}
\begin{itemize}
    \item Frequent words are placed closer to the root, reducing
          the number of computations.
    \item Rare words are placed deeper in the tree, but their
          impact on overall computation is minimal.
    \item The tree minimizes the average path length, making it
          optimal for efficient probability calculations.
\end{itemize}

\section{Word2Vec and Hierarchical Softmax}
\subsection{Tree Depth}
For a vocabulary of size \( N \), a \emph{balanced} binary tree has depth \( \lceil \log_2 N \rceil \):
\begin{itemize}
    \item a vocabulary of 10,000 words gives a depth of 14;
    \item a vocabulary of 100,000 words gives a depth of 17.
\end{itemize}
A Huffman tree is not balanced. Its \emph{average} path length, weighted by word frequencies, lies
between the entropy \( H \) of the unigram distribution and \( H + 1 \) bits, and \( H \leq \log_2 N \): for a
Zipfian vocabulary the average path is noticeably shorter than \( \log_2 N \), which is what makes training
fast. The path of a rare word, however, can be much longer than \( \log_2 N \) (the worst case is \( N - 1 \)),
so ``\( O(\log N) \)'' should be understood as a statement about the average cost.

\subsection{Node Vectors}
Each internal node has a vector of size \( d \) (e.g., 300
in Word2Vec). These vectors are used by binary classifiers
to make decisions (left or right). The total number of
vectors is \( N-1 \), where \( N \) is the vocabulary size.

\subsection{Binary Classifiers}
At each internal node, a binary classifier (e.g., logistic
regression) computes the probability of going left or right:
\[
P(\text{left} \,|\, \text{node}, w_i) = \sigma(\mathbf{v}_{\text{node}}^T
\cdot w_i + b_{\text{node}}),
\]
where:
\begin{itemize}
    \item \( \mathbf{v}_{\text{node}} \) is the node's vector.
    \item \( w_i \) is the context vector.
    \item \( b_{\text{node}} \) is a bias term.
    \item \( \sigma \) is the sigmoid function.
\end{itemize}

\subsection{Training}
During training:
\begin{enumerate}
    \item For a given context vector \( w_i \) and target word \( w \),
          follow the path from the root to the leaf node for \( w \).
    \item Update the parameters of the binary classifiers along the
          path using gradient descent.
    \item Only the classifiers on the path are updated, making
          training efficient.
\end{enumerate}

\subsection{Simultaneous Learning of Tree and Embeddings}
In Word2Vec, the \textbf{Huffman tree} and \textbf{word embeddings} are learned
simultaneously during training. Here's how this works:
\begin{itemize}
    \item \textbf{Tree Initialization}: The Huffman tree is initialized
          based on word frequencies. Frequent words are placed closer
          to the root, and rare words are placed deeper in the tree.
    \item \textbf{Node Vectors}: Each internal node is assigned a vector
          (initialized randomly), which is used by binary classifiers
          to make decisions (left or right).
    \item \textbf{Forward Pass}: For a given context-target pair, the
          model computes the probability of the target word by traversing
          the tree from the root to the leaf node.
    \item \textbf{Backward Pass}: Gradients are propagated back through
          the tree, updating:
          \begin{itemize}
              \item The \textbf{context vector} (input embedding).
              \item The \textbf{node vectors} along the path to the target word.
              \item The \textbf{binary classifiers} at each internal node.
          \end{itemize}
    \item \textbf{Dynamic Adaptation}: Over time, the tree structure,
          node vectors, and word embeddings are jointly optimized to
          minimize the prediction error.
\end{itemize}

\section{Vectors in Word2Vec}
\subsection{What Are the Vectors?}
In Word2Vec, each word is represented as a dense vector in a
continuous vector space. These vectors capture semantic
relationships between words. There are two types of vectors:
\begin{itemize}
    \item \textbf{Input Vectors}: Represent the context words.
    \item \textbf{Output Vectors}: Represent the target words.
\end{itemize}

\subsection{How Are They Learned?}
Word2Vec learns these vectors using one of two algorithms:
\begin{itemize}
    \item \textbf{Continuous Bag of Words (CBOW)}: Predicts a target
          word given its context words.
    \item \textbf{Skip-Gram}: Predicts context words given a target
          word.
\end{itemize}

\subsection{Hierarchical Softmax in Word2Vec}
Hierarchical Softmax is used in both CBOW and Skip-Gram to
efficiently compute the probability of a target word given
its context (or vice versa). It replaces the traditional
Softmax layer with a Huffman tree, reducing computational
complexity.

\section{Algorithms Using Hierarchical Softmax}
Hierarchical Softmax is used in the following algorithms:
\begin{itemize}
    \item \textbf{Word2Vec}: For efficient word embedding learning.
    \item \textbf{Glove}: Although Glove typically uses other
          optimization techniques, variants may use Hierarchical
          Softmax.
    \item \textbf{FastText}: Extends Word2Vec to subword
          representations and uses Hierarchical Softmax for
          large vocabularies.
\end{itemize}

\section{Advantages of Hierarchical Softmax}
\begin{itemize}
    \item \textbf{Efficiency}: Reduces computational complexity
          from \( O(N) \) to \( O(\log N) \).
    \item \textbf{Scalability}: Works well for large vocabularies.
    \item \textbf{Shared Structure}: A single tree is shared across
          all context vectors, reducing memory usage.
\end{itemize}

\section{Conclusion}
Hierarchical Softmax and Huffman trees are powerful tools
for efficiently computing probabilities in large vocabularies.
They are widely used in Word2Vec and other NLP models to
improve computational efficiency and scalability.
\end{notechapter}
